% Generated draft: requires mathematical and source review.
\documentclass[11pt]{article}
\usepackage[margin=25mm]{geometry}
\usepackage{fontspec}
\setmainfont{Latin Modern Roman}
\usepackage{amsmath,amssymb,mathtools,unicode-math}
\setmathfont{Latin Modern Math}
\usepackage{enumitem,needspace,xurl,hyperref}
\setlength{\parindent}{0pt}
\setlength{\parskip}{5pt}
\setlength{\emergencystretch}{3em}
\newcommand{\N}{\mathbb{N}}
\newcommand{\Z}{\mathbb{Z}}
\newcommand{\Q}{\mathbb{Q}}
\newcommand{\R}{\mathbb{R}}
\newcommand{\C}{\mathbb{C}}
\newcommand{\F}{\mathbb{F}}
\newcommand{\A}{\mathbb{A}}
\newcommand{\PP}{\mathbb{P}}
\newcommand{\E}{\mathbb{E}}
\newcommand{\eps}{\varepsilon}
\newcommand{\dd}{\,\mathrm{d}}
\begin{document}
\section*{UnsolvedMath ID: 30000676}
\textbf{Source problem code:} catalogue entry
\subsection{Definitions and mathematical statement}
In the language of groups, quantifier-free formulas are Boolean combinations of equations between group terms, with parameters if permitted. Let $\operatorname{QF\!\text{-}rk}(G)$ denote the quantifier-free rank intended by the source. The extract does not supply its rank recursion, parameter convention or ambient structure, so it is not identified with Morley rank or with an algebraic generator rank.

The first question asks for an analogue of the Tits alternative for finitely generated groups $G$ of finite QF-rank. The standard Tits dichotomy to compare with is
\[
G\text{ is virtually solvable}\quad\text{or}\quad
G\text{ contains a subgroup isomorphic to the nonabelian free group }F_2,
\]
where virtually solvable means having a solvable subgroup of finite index. This displays the usual candidate meaning of ``analogue''; the source does not specify whether it asks for exactly this dichotomy or a modified conclusion.

A group is CSA if every maximal abelian subgroup $A$ is malnormal:
\[
gAg^{-1}\cap A=\{1\}\qquad(g\in G\setminus A).
\]
In the abstract-group convention, a $p$-group is a group whose every element has finite order a power of the prime $p$. The second question asks whether there exists a CSA $p$-group of finite QF-rank. Finitely generated is stated only in the first question and is not carried over to this existence question.

The second question requires a further nontriviality convention: under the displayed CSA definition every abelian group is CSA, and finite cyclic $p$-groups have finite rank under the usual finite-set rank conventions. Thus an intended nonabelian, infinite or other restricted existence problem cannot be inferred from the literal wording. The extract also does not identify a fixed prime or distinguish abstract $p$-groups from pro-$p$ groups. These missing choices are recorded without strengthening the source.
\subsection{Short English statement}
Seek a Tits-type dichotomy for finitely generated finite-QF-rank groups, and ask for CSA $p$-groups of finite QF-rank. The rank and analogue are unspecified, and the second question needs clarification to exclude its elementary abelian examples if that is intended.
\subsection{Sources}
\begin{itemize}
\item Supplied report: \emph{Model Theory and Groups} (2007), Oberwolfach Reports, DOI 10.4171/owr/2007/02, pp. 83--138; proposed year 2007. \url{https://doi.org/10.4171/owr/2007/02}.
\item Local UnsolvedMath registry, permanent ID 30000676, catalogue entry; created 25 August 2026, updated 16 September 2026. Its corrected\_verified note reports comparison with the official report. The supplied 21 August 2026 literature triage labels the extracted questions open and reports no located authoritative complete resolution. No additional external reference is supplied; these dated labels do not resolve the formulation qualifications and have not been independently checked here.
\end{itemize}
\end{document}
